% GENERATED FILE. Edit canonical lessons, then run npm run generate:books.
% canonical-source-hash: fnv1a64:2764c441586f07cc
% canonical-lessons: PA-C262-fis, PA-C262-kiraia, PA-C262-jurmana, PA-C262-rasid, PA-C262-chukauna

\chapter{Fees and Fines}
\label{ch:pa-a2-fees-and-fines}
\hlchaptermodality{262}

\begin{chapteropening}
\textbf{By the end of this chapter:} \emph{I can understand a fee, a fare and a fine.}

Complete the last lesson of chapter 262: I can understand a fee, a fare and a fine.
\end{chapteropening}

\section[\texorpdfstring{\pa{ਫ਼ੀਸ} (fīs)}{ਫ਼ੀਸ (fīs)}]{\pa{ਫ਼ੀਸ} (fīs) — a fee}

\label{lesson:PA-C262-fis}



\begin{quote}
Before the new one: say the Punjabi for to set off, then the Punjabi for departing.
\end{quote}

\subsection*{You'll want to know: \pa{ਫ਼ੀਸ}}
\textbf{\pa{ਫ਼ੀਸ}} — \emph{fīs} — \textquotedblleft{}a fee\textquotedblright{}.

\begin{quote}
\pa{ਲਾਇਬ੍ਰੇਰੀ} \pa{ਦੀ} \pa{ਫ਼ੀਸ} \pa{ਬਹੁਤ} \pa{ਘੱਟ} \pa{ਹੈ।}
\end{quote}

\begin{quote}
\emph{lāibrerī dī fīs bahut ghaṭṭ hai.}
\end{quote}

\begin{quote}
The library fee is very low.
\end{quote}

A feminine noun.

\begin{cousinweb}
\textbf{\pa{ਫ਼ੀਸ}} (\emph{fīs}) is borrowed from English \emph{fees}, the plural of \emph{fee}.
\end{cousinweb}

\begin{grammarlens}[title={What you've built}]
One more word for reading notices and messages.
\end{grammarlens}

\subsection*{Your turn}
\begin{itemize}
\raggedright
  \item \emph{Say it:} \emph{fīs}
  \item \emph{Say it:} \emph{fīs}, once more
  \item \emph{Say it:} \emph{lāibrerī dī fīs bahut ghaṭṭ hai.} — The library fee is very low.
  \item \emph{Recall:} say the Punjabi for to set off, then the Punjabi for departing, then say \emph{fīs} again
\end{itemize}

\begin{tcolorbox}[breakable,colback=teal!4,colframe=teal!35!black,title={Before you move on}]
What does \pa{ਫ਼ੀਸ} mean? (\textquotedblleft{}A fee\textquotedblright{}.) Say it once more.
\end{tcolorbox}
\section[\texorpdfstring{\pa{ਕਿਰਾਇਆ} (kirāiā)}{ਕਿਰਾਇਆ (kirāiā)}]{\pa{ਕਿਰਾਇਆ} (kirāiā) — a fare; rent}

\label{lesson:PA-C262-kiraia}



\begin{quote}
Before the new one: say the Punjabi for departing, then the Punjabi for a fee.
\end{quote}

\subsection*{You'll want to know: \pa{ਕਿਰਾਇਆ}}
\textbf{\pa{ਕਿਰਾਇਆ}} — \emph{kirāiā} — \textquotedblleft{}a fare; rent\textquotedblright{}.

\begin{quote}
\pa{ਸ਼ਹਿਰ} \pa{ਤੱਕ} \pa{ਬੱਸ} \pa{ਦਾ} \pa{ਕਿਰਾਇਆ} \pa{ਸਸਤਾ} \pa{ਹੈ।}
\end{quote}

\begin{quote}
\emph{shahir tak bass dā kirāiā sastā hai.}
\end{quote}

\begin{quote}
The bus fare to town is cheap.
\end{quote}

A masculine noun.

\begin{cousinweb}
\textbf{\pa{ਕਿਰਾਇਆ}} (\emph{kirāiā}) is a loan, usually traced through Persian to Arabic \emph{kirā}, \textquotedblleft{}hire\textquotedblright{}.
\end{cousinweb}

\begin{grammarlens}[title={What you've built}]
One more word for reading notices and messages.
\end{grammarlens}

\subsection*{Your turn}
\begin{itemize}
\raggedright
  \item \emph{Say it:} \emph{kirāiā}
  \item \emph{Say it:} \emph{kirāiā}, once more
  \item \emph{Say it:} \emph{shahir tak bass dā kirāiā sastā hai.} — The bus fare to town is cheap.
  \item \emph{Recall:} say the Punjabi for departing, then the Punjabi for a fee, then say \emph{kirāiā} again
  \item \emph{Write it:} copy \textbf{\pa{ਕਿਰਾਇਆ}} once, checking the whole word
\end{itemize}

\begin{tcolorbox}[breakable,colback=teal!4,colframe=teal!35!black,title={Before you move on}]
What does \pa{ਕਿਰਾਇਆ} mean? (\textquotedblleft{}A fare; rent\textquotedblright{}.) Say it once more.
\end{tcolorbox}
\section[\texorpdfstring{\pa{ਜੁਰਮਾਨਾ} (jurmānā)}{ਜੁਰਮਾਨਾ (jurmānā)}]{\pa{ਜੁਰਮਾਨਾ} (jurmānā) — a fine; a penalty}

\label{lesson:PA-C262-jurmana}



\begin{quote}
Before the new one: say the Punjabi for a fee, then the Punjabi for a fare.
\end{quote}

\subsection*{You'll want to know: \pa{ਜੁਰਮਾਨਾ}}
\textbf{\pa{ਜੁਰਮਾਨਾ}} — \emph{jurmānā} — \textquotedblleft{}a fine; a penalty\textquotedblright{}.

\begin{quote}
\pa{ਜੇ} \pa{ਤੁਹਾਡੇ} \pa{ਕੋਲ} \pa{ਟਿਕਟ} \pa{ਨਹੀਂ} \pa{ਹੈ} \pa{ਤਾਂ} \pa{ਤੁਸੀਂ} \pa{ਜੁਰਮਾਨਾ} \pa{ਭਰੋਗੇ।}
\end{quote}

\begin{quote}
\emph{je tuhāḍe kol ṭikaṭ nahīṁ hai tāṁ tusī̃ jurmānā bharoge.}
\end{quote}

\begin{quote}
If you have no ticket, you will pay a fine.
\end{quote}

A masculine noun.

\begin{cousinweb}
\textbf{\pa{ਜੁਰਮਾਨਾ}} (\emph{jurmānā}) is a loan, usually traced to Persian \emph{jurmāna}, \textquotedblleft{}a fine\textquotedblright{}.
\end{cousinweb}

\begin{grammarlens}[title={What you've built}]
One more word for reading notices and messages.
\end{grammarlens}

\subsection*{Your turn}
\begin{itemize}
\raggedright
  \item \emph{Say it:} \emph{jurmānā}
  \item \emph{Say it:} \emph{jurmānā}, once more
  \item \emph{Say it:} \emph{je tuhāḍe kol ṭikaṭ nahīṁ hai tāṁ tusī̃ jurmānā bharoge.} — If you have no ticket, you will pay a fine.
  \item \emph{Recall:} say the Punjabi for a fee, then the Punjabi for a fare, then say \emph{jurmānā} again
  \item \emph{Write it:} copy \textbf{\pa{ਜੁਰਮਾਨਾ}} once, checking the whole word
\end{itemize}

\begin{tcolorbox}[breakable,colback=teal!4,colframe=teal!35!black,title={Before you move on}]
What does \pa{ਜੁਰਮਾਨਾ} mean? (\textquotedblleft{}A fine; a penalty\textquotedblright{}.) Say it once more.
\end{tcolorbox}
\section[\texorpdfstring{\pa{ਰਸੀਦ} (rasīd)}{ਰਸੀਦ (rasīd)}]{\pa{ਰਸੀਦ} (rasīd) — a receipt}

\label{lesson:PA-C262-rasid}



\begin{quote}
Before the new one: say the Punjabi for a fare, then the Punjabi for a fine.
\end{quote}

\subsection*{You'll want to know: \pa{ਰਸੀਦ}}
\textbf{\pa{ਰਸੀਦ}} — \emph{rasīd} — \textquotedblleft{}a receipt\textquotedblright{}.

\begin{quote}
\pa{ਫ਼ੀਸ} \pa{ਦੀ} \pa{ਰਸੀਦ} \pa{ਮੇਰੀ} \pa{ਜੇਬ} \pa{ਵਿੱਚ} \pa{ਹੈ।}
\end{quote}

\begin{quote}
\emph{fīs dī rasīd merī jeb vich hai.}
\end{quote}

\begin{quote}
The receipt for the fee is in my pocket.
\end{quote}

A feminine noun.

\begin{cousinweb}
\textbf{\pa{ਰਸੀਦ}} (\emph{rasīd}) is a loan, usually traced to Persian \emph{rasīd}, \textquotedblleft{}arrived; a receipt\textquotedblright{}.
\end{cousinweb}

\begin{grammarlens}[title={What you've built}]
One more word for reading notices and messages.
\end{grammarlens}

\subsection*{Your turn}
\begin{itemize}
\raggedright
  \item \emph{Say it:} \emph{rasīd}
  \item \emph{Say it:} \emph{rasīd}, once more
  \item \emph{Say it:} \emph{fīs dī rasīd merī jeb vich hai.} — The receipt for the fee is in my pocket.
  \item \emph{Recall:} say the Punjabi for a fare, then the Punjabi for a fine, then say \emph{rasīd} again
  \item \emph{Write it:} copy \textbf{\pa{ਰਸੀਦ}} once, checking the whole word
\end{itemize}

\begin{tcolorbox}[breakable,colback=teal!4,colframe=teal!35!black,title={Before you move on}]
What does \pa{ਰਸੀਦ} mean? (\textquotedblleft{}A receipt\textquotedblright{}.) Say it once more.
\end{tcolorbox}
\section[\texorpdfstring{\pa{ਚੁਕਾਉਣਾ} (chukāuṇā)}{ਚੁਕਾਉਣਾ (chukāuṇā)}]{\pa{ਚੁਕਾਉਣਾ} (chukāuṇā) — to pay off; to settle a bill}

\label{lesson:PA-C262-chukauna}



\begin{quote}
Before the new one: say the Punjabi for a fine, then the Punjabi for a receipt.
\end{quote}

\subsection*{You'll want to know: \pa{ਚੁਕਾਉਣਾ}}
\textbf{\pa{ਚੁਕਾਉਣਾ}} — \emph{chukāuṇā} — \textquotedblleft{}to pay off; to settle a bill\textquotedblright{}.

\begin{quote}
\pa{ਮੈਂ} \pa{ਹਰ} \pa{ਮਹੀਨੇ} \pa{ਬੈਂਕ} \pa{ਦਾ} \pa{ਕਰਜ਼ਾ} \pa{ਚੁਕਾਉਂਦਾ} \pa{ਹਾਂ।}
\end{quote}

\begin{quote}
\emph{maiṁ har mahīne baiṅk dā karzā chukāundā hāṁ.}
\end{quote}

\begin{quote}
I pay off the bank debt every month.
\end{quote}

Past (I …): \pa{ਮੈਂ} \pa{ਜੁਰਮਾਨਾ} \pa{ਚੁਕਾਇਆ}, \emph{maiṁ jurmānā chukāiā}.

\begin{cousinweb}
\textbf{\pa{ਚੁਕਾਉਣਾ}} (\emph{chukāuṇā}) is inherited from a reconstructed \emph{*cukk}, \textquotedblleft{}fall short of, stop\textquotedblright{}; CDIAL lists Punjabi \emph{cukāuṇā}, \textquotedblleft{}to finish with, settle\textquotedblright{}, under it (CDIAL 4848).
\end{cousinweb}

\begin{grammarlens}[title={What you've built}]
One more word for reading notices and messages.
\end{grammarlens}

\subsection*{Your turn}
\begin{itemize}
\raggedright
  \item \emph{Say it:} \emph{chukāuṇā}
  \item \emph{Say it:} \emph{chukāuṇā}, once more
  \item \emph{Say it:} \emph{maiṁ har mahīne baiṅk dā karzā chukāundā hāṁ.} — I pay off the bank debt every month.
  \item \emph{Recall:} say the Punjabi for a fine, then the Punjabi for a receipt, then say \emph{chukāuṇā} again
  \item \emph{Write it:} copy \textbf{\pa{ਚੁਕਾਉਣਾ}} once, checking the whole word
\end{itemize}

\begin{tcolorbox}[breakable,colback=teal!4,colframe=teal!35!black,title={Before you move on}]
What does \pa{ਚੁਕਾਉਣਾ} mean? (\textquotedblleft{}To pay off; to settle a bill\textquotedblright{}.) Say it once more.
\end{tcolorbox}
